% ID: SEL-D0045-OB-E02
% ESTADO: validada. S046, cotejada con examen y criterios.
El dominio es $\mathbb R$. Al factorizar,
$f(x)=x[(x-3)^2+3]$, así que el único corte con ambos ejes
es el origen. También $f(x)=(x-2)^3+8$.

Como $f'(x)=3(x-2)^2$ es positiva salvo en $2$ y no cambia de
signo, $f$ crece estrictamente en todo $\mathbb R$ y carece
de extremos relativos. La segunda derivada es $6(x-2)$.
La curva es cóncava hacia abajo antes de $2$ y convexa
después. El punto de inflexión es $\boxed{(2,8)}$, con
tangente horizontal.

La gráfica es la cúbica básica trasladada dos unidades
a la derecha y ocho hacia arriba:
\begin{center}
\setlength{\unitlength}{7mm}
\begin{picture}(8,8)(-2,-2)
\qbezier(-1.5,0)(1.75,0)(5,0)\put(5.1,0){$x$}
\qbezier(0,-1.5)(0,2)(0,5.5)\put(0,5.6){$f$}
\qbezier(-0.25,-0.84766)(-0.125,-0.37402)(0,0.00000)
\qbezier(0,0.00000)(0.25,0.74219)(0.5,1.15625)
\qbezier(0.5,1.15625)(0.75,1.57031)(1,1.75000)
\qbezier(1,1.75000)(1.25,1.92969)(1.5,1.96875)
\qbezier(1.5,1.96875)(1.75,2.00781)(2,2.00000)
\qbezier(2,2.00000)(2.25,1.99219)(2.5,2.03125)
\qbezier(2.5,2.03125)(2.75,2.07031)(3,2.25000)
\qbezier(3,2.25000)(3.25,2.42969)(3.5,2.84375)
\qbezier(3.5,2.84375)(3.75,3.25781)(4,4.00000)
\qbezier(4,4.00000)(4.15,4.44831)(4.3,5.04175)
\put(2,2.2){\makebox(0,0)[b]{$(2,8)$}}
\put(0.1,-0.2){$(0,0)$}
\end{picture}
\end{center}
